\chapter{Fifth Disease}
\index{Fifth Disease}

\section*{Definition and Overview}
Fifth disease, also called slapped cheek disease or erythema infectiosum, is a common, mild viral illness of childhood caused by parvovirus B19. It is called "fifth disease" because it was the fifth of the classic childhood rash illnesses historically described. The illness is most recognisable by its bright red rash on the cheeks, which looks as though the child has been slapped, followed by a lace-like rash on the body and limbs. In most children, fifth disease is so mild that no treatment is needed beyond comfort measures. However, parvovirus B19 can cause problems in certain groups: it can trigger an acute drop in red blood cells (aplastic crisis) in people with blood disorders, and infection during pregnancy can affect the unborn baby. This chapter explains the illness, its complications, and who needs special care.

\section*{Epidemiology}
Fifth disease is common worldwide and occurs mainly in school-aged children, with most cases in late winter and spring. Outbreaks are frequent in schools and childcare, and by adulthood, the majority of people have been infected and are immune. Around 1--5\% of pregnant women are infected during pregnancy, although most will have no effect on the baby. Parvovirus B19 causes significant concern in pregnancy and in people with sickle cell disease and other haemolytic disorders, where the virus can precipitate a severe drop in haemoglobin. Aplastic crisis is most common in people with chronic haemolytic anaemia, including sickle cell disease and hereditary spherocytosis.

\section*{Aetiology and Pathophysiology}
\begin{itemize}
    \item \textbf{The virus:} parvovirus B19 infects only humans and targets rapidly dividing cells, particularly the red blood cell precursors in the bone marrow
    \item \textbf{Transmission:} spread by respiratory droplets from coughing and sneezing, and by contact with infected secretions; it is contagious before the rash appears
    \item \textbf{Incubation:} symptoms appear 4--14 days after exposure
    \item \textbf{The rash:} the rash is an immune response that appears after the virus has been cleared, so a child with the rash is usually no longer contagious
    \item \textbf{Effects on red cells:} the virus temporarily stops red cell production; in healthy people this is unnoticed, but in those with haemolytic anaemia it causes a sudden, severe drop in haemoglobin (aplastic crisis)
    \item \textbf{Pregnancy effects:} the virus can cross the placenta; infection before 20 weeks carries a small risk of severe fetal anaemia, hydrops, or miscarriage
    \item \textbf{Immunity:} infection produces lifelong immunity, and most adults are already immune
\end{itemize}

\section*{Clinical Features}
\begin{itemize}
    \item Mild fever, headache, and a runny nose in the early days, sometimes with no symptoms at all
    \item The classic "slapped cheek" rash: bright red, flat redness on both cheeks, lasting a few days
    \item A lace-like or blotchy rash on the arms, trunk, and thighs, which may come and go for 1--3 weeks and can be aggravated by heat or exercise
    \item Joint pain and swelling, mainly in adult women, affecting the hands, wrists, and knees, which can last weeks
    \item In most children: a mild illness that resolves without treatment
    \item In people with haemolytic anaemia: sudden pallor, fatigue, breathlessness, and a racing heart from aplastic crisis
    \item In pregnancy: often asymptomatic or mild in the mother, with the risk mainly to the fetus
\end{itemize}

\section*{Differential Diagnosis}
The rash of fifth disease can resemble other childhood rashes.
\begin{itemize}
    \item Measles, rubella, and scarlet fever
    \item Roseola infantum
    \item Viral exanthems and allergic rashes
    \item Hand, foot and mouth disease
    \item In adults with joint symptoms: rheumatoid arthritis and other inflammatory arthritis
    \item In pregnancy: other infections such as rubella, cytomegalovirus, and toxoplasmosis
\end{itemize}

\section*{When to Seek Medical Care}
Most children with fifth disease need no medical care. Medical review is needed for people with blood disorders, women who are pregnant or may be pregnant, and anyone with symptoms of anaemia such as marked pallor, fatigue, or breathlessness.

\textbf{Contact your local emergency services for:}
\begin{itemize}
    \item Sudden severe pallor, breathlessness, fainting, or a racing heart in a person with a blood disorder
    \item A pregnant woman with possible parvovirus exposure who develops severe symptoms
    \item Chest pain or collapse
    \item Any emergency symptoms
\end{itemize}

\section*{Diagnosis and Investigations}
\begin{itemize}
    \item \textbf{Clinical diagnosis:} the characteristic slapped cheek rash in a child is usually sufficient, and no testing is needed
    \item \textbf{Blood tests:} parvovirus B19 antibody tests (IgM and IgG) and PCR in pregnant women, people with blood disorders, and adults with joint symptoms
    \item \textbf{Full blood count and reticulocyte count:} in people with suspected aplastic crisis, showing a sudden drop in red cells and low reticulocytes
    \item \textbf{Ultrasound and fetal monitoring:} in pregnancy, serial scans and measurement of fetal blood flow to detect fetal anaemia
    \item \textbf{Serology for other infections:} in pregnancy, to exclude rubella and other causes of fetal harm
\end{itemize}

\section*{Management}
\begin{itemize}
    \item \textbf{Children:} no specific treatment; rest, fluids, and paracetamol for discomfort if needed
    \item \textbf{Joint symptoms:} rest and anti-inflammatory medicines for adults with arthritis; symptoms usually resolve within weeks
    \item \textbf{Aplastic crisis:} hospital care with blood transfusion for people with haemolytic anaemia whose red cell count falls dangerously
    \item \textbf{Pregnancy:} close monitoring by an obstetrician, with serial ultrasound; if fetal anaemia develops, intrauterine transfusion can be life-saving
    \item \textbf{Infection control:} good hand hygiene, and exclusion from school or childcare until the fever settles, since the contagious phase is before the rash
    \item \textbf{Immunocompromised people:} chronic parvovirus infection can persist and may need specialist treatment with immunoglobulin
\end{itemize}

\section*{Complications}
\begin{itemize}
    \item Aplastic crisis in people with sickle cell disease, hereditary spherocytosis, and other haemolytic disorders
    \item Chronic anaemia in people with weakened immune systems
    \item Fetal anaemia, hydrops fetalis, and miscarriage, in a minority of infections acquired before 20 weeks of pregnancy
    \item Joint pain lasting weeks to months in some adults
    \item Very rarely, myocarditis and other organ involvement
\end{itemize}

\section*{Prognosis and Outcomes}
The outlook for fifth disease is excellent. Healthy children recover fully with no treatment, and the illness confers lifelong immunity. In people with blood disorders, aplastic crisis resolves once the bone marrow recovers, with transfusion support. In pregnancy, the great majority of infections cause no harm to the baby; for the small minority in which fetal anaemia develops, specialist monitoring and transfusion rescue the baby in most cases. With appropriate care for the at-risk groups, the complications of parvovirus B19 are largely preventable or treatable.

\section*{Prevention and Self-Care}
\begin{itemize}
    \item Wash hands frequently, especially in childcare and school settings
    \item Cover coughs and sneezes, and keep children home while unwell with fever
    \item Pregnant women who work with children should be aware of the infection and discuss any exposure with a doctor
    \item People with blood disorders should seek prompt medical review after known exposure
    \item Do not use aspirin in children with viral illness
    \item Provide fluids, rest, and comfort measures for the affected child
    \item Follow the advice of the obstetric team if infection occurs in pregnancy
\end{itemize}

\section*{Sources and Further Reading}
The following reputable sources provide further information for healthcare students and patients seeking reliable, current advice about fifth disease.
\begin{itemize}
    \item Healthdirect Australia. \textit{Fifth disease.} https://www.healthdirect.gov.au/fifth-disease
    \item The Royal Children's Hospital Melbourne. \textit{Slapped cheek disease fact sheet.} https://www.rch.org.au/
    \item NSW Health. \textit{Parvovirus B19 fact sheet.} https://www.health.nsw.gov.au/
    \item NHS. \textit{Slapped cheek syndrome.} https://www.nhs.uk/conditions/slapped-cheek-syndrome/
    \item Mayo Clinic. \textit{Parvovirus infection.} https://www.mayoclinic.org/
    \item Centers for Disease Control and Prevention. \textit{Parvovirus B19.} https://www.cdc.gov/
\end{itemize}
